\chapter{Conclusiones y trabajos futuros}
\label{chap:conclusiones-trabajos-futuros}

En este capítulo se presenta un breve resumen de la consecución de los requisitos que se establecieron para el proyecto y una reflexión acerca del trabajo realizado.


\section{Consecución de objetivos}
\label{sec:consecucion-objetivos}

Mediante los resultados presentados en el capítulo \textit{\ref{chap:validacion-resultados} \nameref{chap:validacion-resultados}}, se puede comprobar que los objetivos marcados al inicio del proyecto han sido alcanzados, obteniendo una aplicación que sirve como base para una futura potente herramienta para el análisis y simulación de control aéreo.

Así, se ha implementado la navegación autónoma lateral de las aeronaves en base a puntos de ruta y procedimientos estandarizados de aproximación (STAR) y salida (SID), navegación básica en base a rumbos, altitudes y velocidades, así como una amplia y sólida base para albergar todos los elementos que componen la navegación aérea como fijos, radioayudas, procedimientos RNAV, etc.

La navegación lateral se realiza en base al cálculo constante de la rotación que debería tener la aeronave en ese momento para estar mirando hacia el siguiente punto de ruta. De este modo, se establece un rumbo objetivo que se va corrigiendo constantemente a medida que la aeronave realiza virajes, y de tal forma que siga siendo válida si se implementa viento en el simulador, ya que la función de navegación lateral irá corrigiendo la deriva constantemente. En este caso se apreciaría una navegación con variación de rumbo continuamente.

De un modo similar, durante la navegación lateral autónoma la aeronave vuela secuencialmente los diferentes puntos de ruta de un procedimiento estándar mediante la comprobación periódica de la distancia horizontal hacia el siguiente punto. Cuando dicha distancia es inferior a \SI{1}{NM} se da por sobrevolado el punto, eliminándose de la lista de puntos siguientes, pasando a ser el siguiente punto de ruta al nuevo punto activo.

La navegación vertical se ha modelado en base la velocidad vertical de cada tipo de aeronave, aceptando el modificador \textit{expedite} para agilizar la maniobra en caso necesario. Así, se establece una tasa de ascenso o descenso estándar según la categoría y se aplica un ligero componente aleatorio en cada periodo de actualización radar para dar un poco de realismo y no sean números fijos.

Del mismo modo, el incremento o reducción de velocidad y los virajes se han modelado en base a las prestaciones de cada categoría de aeronave, añadiendo cierta variabilidad aleatoria.

En el apartado de representación gráfica de los elementos de ruta y aeroportuarios se ha conseguido un gran nivel de fidelidad, siendo las distancias y rumbos acorde a las cartas aeronáuticas oficiales, lo que le confiere un alto nivel de calidad a los datos de trayectorias y tiempos de vuelo que puedan ser tomados para futuros análisis.


En cuanto a la interacción del usuario con la aplicación, se ha cuidado desde el primer momento, prestando especial atención en la sencillez y rapidez para establecer comandos a la aeronave, sin descuidar el apartado gráfico. La principal premisa, aparte de las que respectan a los aspectos técnicos, ha sido crear un entorno agradable para el usuario, de tal modo que la aplicación en su conjunto estuviese en armonía y no hubiese elementos disonantes.


\clearpage

\section{Trabajos futuros}
\label{sec:trabajos-futuros}

Debido a la envergadura de este proyecto, existen ciertas ideas que no tenían cabida aquí, pero que pueden presentarse como posibles mejoras de las funcionalidades de la aplicación. A continuación, se presentan las identificadas como más relevantes en el mismo formato que se hicieron los requisitos.


\newcounter{TF_index}
\setcounter{TF_index}{1}

% Encabezado Trabajos Futuros
\begin{table}[!hbtp]
	\footnotesize
	\setlength{\arrayrulewidth}{0.4mm}
	%\setlength{\tabcolsep}{5pt}
	\renewcommand{\arraystretch}{1.2}
	\begin{center}
		\caption{Tabla con los Trabajos Futuros.}
		\label{tab:Fase1}
		\begin{tabular}{|c|}
			%% Title
			\hline 
			%\multicolumn{2}{|c|}{\cellcolor{gray!30}{\textbf{TRABAJOS FUTUROS}}}
			\hspace{4.4cm}
			{\cellcolor{gray!30}{\textbf{TRABAJOS FUTUROS}}}
			\hspace{4.4cm}
			\\ 
			\hline
		\end{tabular}
	\end{center}
\end{table}

\begin{table}[!hbtp]
	\footnotesize
	\setlength{\arrayrulewidth}{0.4mm}
	%\setlength{\tabcolsep}{5pt}
	\renewcommand{\arraystretch}{1.2}			
	\begin{center}
		%\caption{.}
		%\label{tab:}
		\begin{tabular}{|c|p{0.7\textwidth}|}			
			\hline
			\multicolumn{2}{|l|}{\cellcolor{gray!10}{\textbf{ID: TF.\arabic{TF_index}}}} \stepcounter{TF_index} 
			\\ 
			\hline		
			\cellcolor{gray!10}{\textbf{Título}} 
			& Patrones de espera 
			\\
			\cellcolor{gray!10}{\textbf{Desc.}} 
			& Implementar patrones de espera estándar sobre puntos de ruta susceptibles de ello (FIX, VOR...). Deben tener una duración aproximada de cuatro minutos, uno por cada segmento del hipódromo, con virajes de tres grados por segundo. 
			\\
			\hline		
		\end{tabular}
	\end{center}
\end{table}	

\begin{table}[!hbtp]
	\footnotesize
	\setlength{\arrayrulewidth}{0.4mm}
	%\setlength{\tabcolsep}{5pt}
	\renewcommand{\arraystretch}{1.2}			
	\begin{center}
		%\caption{.}
		%\label{tab:}
		\begin{tabular}{|c|p{0.7\textwidth}|}			
			\hline
			\multicolumn{2}{|l|}{\cellcolor{gray!10}{\textbf{ID: TF.\arabic{TF_index}}}} \stepcounter{TF_index} 
			\\ 
			\hline		
			\cellcolor{gray!10}{\textbf{Título}} 
			& Declinación magnética.  
			\\
			\cellcolor{gray!10}{\textbf{Desc.}} 
			& Implementar la declinación magnética como un desfase (\textit{offset}) respecto a la orientación mágnetica un valor definido para el aeropuerto del escenario.  
			\\
			\hline		
		\end{tabular}
	\end{center}
\end{table}	

\begin{table}[!hbtp]
	\footnotesize
	\setlength{\arrayrulewidth}{0.4mm}
	%\setlength{\tabcolsep}{5pt}
	\renewcommand{\arraystretch}{1.2}			
	\begin{center}
		%\caption{.}
		%\label{tab:}
		\begin{tabular}{|c|p{0.7\textwidth}|}			
			\hline
			\multicolumn{2}{|l|}{\cellcolor{gray!10}{\textbf{ID: TF.\arabic{TF_index}}}} \stepcounter{TF_index} 
			\\ 
			\hline		
			\cellcolor{gray!10}{\textbf{Título}} 
			& Modelar meteorología. 
			\\
			\cellcolor{gray!10}{\textbf{Desc.}} 
			& Modelar el viento en el sector de control, de tal forma que el rumbo (\textit{heading}) y la derrota (\textit{track}) diferirán al componerse los vectores velocidad.

Puede implementarse en varias fases:

\begin{itemize}
	\item Viento fijo en una sola capa: el viento se mantiene invariante durante toda la simulación y afecta por igual a todas las altitudes.

	\item Viento fijo en tres capas: el viento se mantiene invariante pero difiere en tres rangos de altitud diferente.

	\item Viento según meteorología real descargado desde algún API Web como el de la agencia NOAA \cite{NOAA}.
\end{itemize} 
			\\
			\hline		
		\end{tabular}
	\end{center}
\end{table}	

\begin{table}[!hbtp]
	\footnotesize
	\setlength{\arrayrulewidth}{0.4mm}
	%\setlength{\tabcolsep}{5pt}
	\renewcommand{\arraystretch}{1.2}			
	\begin{center}
		%\caption{.}
		%\label{tab:}
		\begin{tabular}{|c|p{0.7\textwidth}|}			
			\hline
			\multicolumn{2}{|l|}{\cellcolor{gray!10}{\textbf{ID: TF.\arabic{TF_index}}}} \stepcounter{TF_index} 
			\\ 
			\hline		
			\cellcolor{gray!10}{\textbf{Título}} 
			& Locuciones de texto, TTS (\textit{Text-To-Speech}).  
			\\
			\cellcolor{gray!10}{\textbf{Desc.}} 
			& Integración de locuciones de un motor TTS (\textit{Text-To-Speech}) al que se le pase como parámetros los textos ya preparados.  
			\\
			\hline		
		\end{tabular}
	\end{center}
\end{table}	

\begin{table}[!hbtp]
	\footnotesize
	\setlength{\arrayrulewidth}{0.4mm}
	%\setlength{\tabcolsep}{5pt}
	\renewcommand{\arraystretch}{1.2}			
	\begin{center}
		%\caption{.}
		%\label{tab:}
		\begin{tabular}{|c|p{0.7\textwidth}|}			
			\hline
			\multicolumn{2}{|l|}{\cellcolor{gray!10}{\textbf{ID: TF.\arabic{TF_index}}}} \stepcounter{TF_index} 
			\\ 
			\hline		
			\cellcolor{gray!10}{\textbf{Título}} 
			& Modelar masa de las aeronaves. 
			\\
			\cellcolor{gray!10}{\textbf{Desc.}} 
			& Lo ideal sería establecer una ruta origen – destino a cada aeronave y según la distancia calcular el combustible necesario en función del tipo de aeronave. De esta forma, establecer la masa actual de la misma. 

Para simplificar, en este punto se puede tomar como referencia el MTOW (\textit{Maximum Take-Off Weight}) y el ZFW (\textit{Zero Fuel Weight}) y asignarla la masa actual de la aeroanve de la siguiente manera:
	\begin{itemize}
\item Aeronaves en llegada con combustible entre MTOW y ZFW, por ejemplo, cercano al 20\% de capacidad de combustible.
\item Aeronaves en salida con combustible entre MTOW y ZFW, por ejemplo, cercano al 80-90\% de MTOW.
	\end{itemize}
			\\
			\hline		
		\end{tabular}
	\end{center}
\end{table}	

\begin{table}[!hbtp]
	\footnotesize
	\setlength{\arrayrulewidth}{0.4mm}
	%\setlength{\tabcolsep}{5pt}
	\renewcommand{\arraystretch}{1.2}			
	\begin{center}
		%\caption{.}
		%\label{tab:}
		\begin{tabular}{|c|p{0.7\textwidth}|}			
			\hline
			\multicolumn{2}{|l|}{\cellcolor{gray!10}{\textbf{ID: TF.\arabic{TF_index}}}} \stepcounter{TF_index} 
			\\ 
			\hline		
			\cellcolor{gray!10}{\textbf{Título}} 
			& Modelar velocidad mínima de vuelo en base a aerodinámica. 
			\\
			\cellcolor{gray!10}{\textbf{Desc.}} 
			& Con la masa asignada a cada aeronave se puede calcular la sustentación mínima necesaria para vuelo recto y nivelado y, a partir de los datos de cada tipo de aeronave, obtener unos coeficientes de las superficies sustentadoras y estimar una velocidad mínima de vuelo. Como referencia, puede consultarse la base de datos de prestaciones de aeronaves, BADA (\textit{Base of Aircraft Data}) \cite{BADA}. 
			\\
			\hline		
		\end{tabular}
	\end{center}
\end{table}	

\begin{table}[!hbtp]
	\footnotesize
	\setlength{\arrayrulewidth}{0.4mm}
	%\setlength{\tabcolsep}{5pt}
	\renewcommand{\arraystretch}{1.2}			
	\begin{center}
		%\caption{.}
		%\label{tab:}
		\begin{tabular}{|c|p{0.7\textwidth}|}			
			\hline
			\multicolumn{2}{|l|}{\cellcolor{gray!10}{\textbf{ID: TF.\arabic{TF_index}}}} \stepcounter{TF_index} 
			\\ 
			\hline		
			\cellcolor{gray!10}{\textbf{Título}} 
			& Sectorización y zonas D, P, R. 
			\\
			\cellcolor{gray!10}{\textbf{Desc.}} 
			& Implementar dibujado y restricciones de altitud en zonas de altitud mínima y de consideración especial D \textit{Dangerous}, P \textit{Prohibited}, R \textit{Restricted}. 
			\\
			\hline		
		\end{tabular}
	\end{center}
\end{table}	

\begin{table}[!hbtp]
	\footnotesize
	\setlength{\arrayrulewidth}{0.4mm}
	%\setlength{\tabcolsep}{5pt}
	\renewcommand{\arraystretch}{1.2}			
	\begin{center}
		%\caption{.}
		%\label{tab:}
		\begin{tabular}{|c|p{0.7\textwidth}|}			
			\hline
			\multicolumn{2}{|l|}{\cellcolor{gray!10}{\textbf{ID: TF.\arabic{TF_index}}}} \stepcounter{TF_index} 
			\\ 
			\hline		
			\cellcolor{gray!10}{\textbf{Título}} 
			& Importar escenarios con posibilidad de modificación. 
			\\
			\cellcolor{gray!10}{\textbf{Desc.}} 
			& Las aeronaves, compañías, aeropuertos, pistas, SID, STAR, radioayudas y todos los elementos dependientes del escenario a simular pueden implementarse en ficheros con extensión \texttt{.json} para que sean accesibles y modificables por el usuario. 
			\\
			\hline		
		\end{tabular}
	\end{center}
\end{table}	

\begin{table}[!hbtp]
	\footnotesize
	\setlength{\arrayrulewidth}{0.4mm}
	%\setlength{\tabcolsep}{5pt}
	\renewcommand{\arraystretch}{1.2}			
	\begin{center}
		%\caption{.}
		%\label{tab:}
		\begin{tabular}{|c|p{0.7\textwidth}|}			
			\hline
			\multicolumn{2}{|l|}{\cellcolor{gray!10}{\textbf{ID: TF.\arabic{TF_index}}}} \stepcounter{TF_index} 
			\\ 
			\hline		
			\cellcolor{gray!10}{\textbf{Título}} 
			& Exportar datos a ficheros externos. 
			\\
			\cellcolor{gray!10}{\textbf{Desc.}} 
			& Implementar el guardado a ficheros externos del tipo a convenir, generalmente de valores separados por comas (\texttt{.csv}), de todos los datos generados por las aeronaves en vuelo: distancia recorrida, tiempo en vuelo, puntos de ruta navegado, etc.  
			\\
			\hline		
		\end{tabular}
	\end{center}
\end{table}	








%\begin{table}[!hbtp]
%	\footnotesize
%	\setlength{\arrayrulewidth}{0.4mm}
%	%\setlength{\tabcolsep}{5pt}
%	\renewcommand{\arraystretch}{1.2}			
%	\begin{center}
%		%\caption{.}
%		%\label{tab:}
%		\begin{tabular}{|c|p{0.7\textwidth}|}			
%			\hline
%			\multicolumn{2}{|l|}{\cellcolor{gray!10}{\textbf{ID: TF.\arabic{TF_index}}}} \stepcounter{TF_index} 
%			\\ 
%			\hline		
%			\cellcolor{gray!10}{\textbf{Título}} 
%			&  
%			\\
%			\cellcolor{gray!10}{\textbf{Desc.}} 
%			&  
%			\\
%			\hline		
%		\end{tabular}
%	\end{center}
%\end{table}	

\clearpage
